\documentclass[leqno]{jsarticle}
\usepackage{amssymb}
\usepackage{amsthm}
\usepackage{amsmath}
\usepackage{comment}
%\usepackage{showkeys}
	%可換図式用
		%\usepackage[all]{xy}
	%重くなるので使わいときはコメント化しておく。
%定理環境 thmはあまり使わずpropをなるべく使う。
%主張は基本的に証明の中で使い、そのため番号を付けない。
%注意環境を付けてみた。
\newtheorem{thm}{定理}
\newtheorem{prop}[thm]{命題}
\newtheorem{cor}[thm]{系}
\newtheorem{lem}[thm]{補題}
\newtheorem{df}[thm]{定義}
\newtheorem{exam}[thm]{例}
\newtheorem{fact}[thm]{事実}
\newtheorem*{claim}{主張}
\newtheorem*{rmk}{注意}
\renewcommand{\proofname}{{\bf 証明}}
%矢印たち。
%\toは\rightarrowと同じ。\getsは\leftarrowと同じ。同値は\iff
%上向き矢はuparrow逆はdownarrow
\newcommand{\To}{\Longrightarrow}
\newcommand{\Gets}{\Longleftarrow}
%黒板太字
\newcommand{\nn}{\mathbb{N}}
\newcommand{\zz}{\mathbb{Z}}
\newcommand{\qq}{\mathbb{Q}}
\newcommand{\rr}{\mathbb{R}}
\newcommand{\cc}{\mathbb{C}}
%無限大
\newcommand{\mgn}{\infty}
%ギリシャン
\newcommand{\dpst}{\displaystyle}
\newcommand{\ep}{\varepsilon}
\newcommand{\al}{\alpha}
\newcommand{\be}{\beta}
\newcommand{\de}{\delta}
\newcommand{\la}{\lambda}
\newcommand{\La}{\Lambda}
\newcommand{\ga}{\gamma}
\newcommand{\lil}{\la\in \La}
%商集合
\newcommand{\qt}[2]{#1/\! #2}
%enumerate環境
\renewcommand{\labelenumi}{{\rm (\arabic{enumi})}}
%以降alignじゃなくてenumerate を使え。
%なんか演算
\DeclareMathOperator{\card}{card}
\DeclareMathOperator{\supp}{supp}
%なんか演算２
\DeclareMathOperator{\bd}{Bdry}
\DeclareMathOperator{\cl}{CL}
\DeclareMathOperator{\nai}{Int}
\DeclareMathOperator{\di}{diam}
\DeclareMathOperator{\re}{Re}
\DeclareMathOperator{\im}{Im}
%差集合は\setminusで統一する。等号の否定は\neq
%真部分集合は\subsetneqq　偏微分は\partial
%ディスプレイスタイル。\dfracでディスプレイ分数
%空集合は\Oを使う！！！！！！！！！！！！

\title{クラークソンの不等式(未完成)}
\author{電波通信 ナンブ キトラ }
\date{\today}
			\begin{document}
\maketitle

この文書は，
$L^{p}$空間の一様凸性の証明に用いられるクラークソンの不等式を，
三線定理を用いて証明することを目的とする．


%%%
\section{三線定理とか}
この章の話
\cite{miya}を大いに参考にした．
この章で領域といったら
$\cc$
の連結な開集合のことを指すことにする．
まずは基本的な複素関数論をする．
正則関数といったら複素関数論の正則関数とする．
\begin{prop}
領域
$D$
上で定義された
正則関数
$f$
について，
$|f|$
が
$D$
上で最大値をとるならば，
$f$
は定数関数である．
\end{prop}
%%proof
\begin{proof}
$|f|$ 
は点$a\in D$で最大値を達成するとする．
この時十分小さい$r>0$
を取り，
$B(a, r)\subset D$とする．
この時
$B(a,r)$
上で
$f$
のテイラー展開は収束するので，
$B(a,r)$
上で以下のような等式が成り立つ．
\[
f(z)=\sum_{i=0}c_i(z-a)^i. 
\]
これを用いると，$z$を
($a$を中心とする)
極座標で書いたときに
\[
|f(a+re^{i\theta})|^2=\sum_{n,m}c_n\bar{c}_me^{i(n-m)\theta}r^{n+m}
\]
が成り立つ．
よって
\[
\frac{1}{2\pi}\int_{0}^{2\pi}|f(a+re^{i\theta})|^2\ d\theta=\sum_{i}|c_i|^2r^{2i}
\]
である．
ここで
\[
\frac{1}{2\pi}\int_{0}^{2\pi}|f(a+re^{i\theta})|^2\ d\theta\le |f(a)|=|c_0|^2
\]
なので
\[
\sum_{i}|c_i|^2r^{2i}\le |c_0|^2
\]
である．
よって
$i\ge 1$
ならば
$c_{1}=0$
である．
故に$f$は局所定数であることがわかり，
正則関数の剛性から
$f$
が定数であることがわかる．
\end{proof}


上記の最大値原理から以下の系を得る．
%%%%
\begin{cor}
$D$
を
$\cc$
内の領域とし，
$f$
は
$\cl(D)$
上の定数でない連続関数で，
$D$
上で正則であるとする．
このとき
$f$
が
$\cl(D)$
上で最大値をとるならば，
$\partial D$上で最大値をとる．
\end{cor}
%%%


次の定理がこの章の目的である三線定理と呼ばれるものである．
\begin{thm}[三線定理]
$a$, 
$b$を実数とする．
そして
$S=\{z\in \cc\mid a\le \re z\le b \}$とする．
さらに
$f$
を
$S$
上の有界な連続関数で，
$\nai (S)$上で正則な関数とする．
また，
$a\le c\le b$
となる
$c$
について
\[
M(c)=\sup_y|f(c+iy)|<\infty
\]
とする．
このとき
\[
M((1-t)a+tb)\le M(a)^{1-t}M(b)^{t}
\]
が成り立つ．
\end{thm}
\begin{proof}
適当にアフィン変換することにより，
$a=0$, $b=1$としても良い．

まず最初に$M(0)$と$M(1)$が共に正の時に定理を示す．
このとき
\[
f(z)M(0)^{-(1-z)}M(1)^{-z}
\]
を考えることによって、$M(0)=M(1)=1$としても良い
（正の実数$R$に対して$|R^z|=R^{\re(z)}$なので）．

ここで
\[
f_n(z)=f(z)\exp(z(z-1)/n)
\]
とする．
このとき$z(z-1)=x(x-1)-y^2+i(2xy-1)$
なので
\[
|\exp(z(z-1)/n)|\le \exp(x(x-1))\exp(-y^2)\le \exp(-y^2)
\]
である．
よって
$f_n$は $\re z=0,1$を満たす$z$について$|f_n(z)|\le 1$を満たす．
さらに
\[
\lim_{R\to \infty}\sup_{|\im z|=R}|f_n(z)|=0
\]
が成り立つ．
よって，
$|f_n(z)|$
は
$S$
上で最大値を持つ．
最大値原理より，それは境界上で最大値をとるので，
任意の
$z\in S$
について
$|f_n(z)|\le 1$
である．
ここで，
点$z$を固定するごとに
$f_n(z)\to f(z)\ (n\to \infty)$が成り立つので，
任意の$z\in S$について$|f(z)|\le 1$である．
これが求める不等式である．

次に
$M(0)\times M(1)=0$
の時を考察しよう．
数
$\ep>0$
を十分小さくとると
(具体的には，$M(1)$,$M(0)$がともにはゼロでない時には$\ep<\max\{M(0),M(1)\}$
と取れば良い．両方ともゼロなら特に制限はない．)，
$f+\ep$について
$M(0)$
と
$M(1)$
が正になる．
そして$\ep\to 0$とすれば求める不等式が出てくる．
\end{proof}

測度空間
$(X, \mu)$
について，
$L^0(X, \mu)$
で
$(X, \mu)$
上の可測関数を表す．
\begin{thm}[Riesz--Thorinの定理]
$(X,\mathcal{A},\mu)$と$(Y,\mathcal{B}, \nu)$
を測度空間とする。
そして
$p_{0}, q_{0}, p_{1},q_{1} \in [1,\infty] $で
$\theta\in [0,1]$とする。
この時
\[
\frac{1}{p}=\frac{1-\theta}{p_0}+\frac{\theta}{p_1}
\]
と
\[
\frac{1}{q}=\frac{1-\theta}{q_0}+\frac{\theta}{q_1}
\]
が成り立つ
とする．
このとき線形写像
\[
T\colon L^0(X,\mu)\to L^0(Y,\nu)
\]
と
$(X,\mathcal{A}, \mu)$
の単関数
$f$
と
$(Y,\mathcal{B}, \nu)$
の単関数
$g$
について，
\[
\left|\int_{Y}gTf\ d\nu \right|
\le \|T\|_{p_0, q_0}^{1-\theta}\|T\|_{p_1.q_1}^{\theta}\|f\|_p\|g\|_r
\]
が成り立つ．
ここで$r$は$q$の共役指数である．
\end{thm}
\begin{proof}
まず
$\|f\|_p=\|g\|_r=1$
としても一般性を失わない．
そして
\[
f(x)=\sum_{j=1}^na_j\chi_{A_i}(x)
\]
とおく．
ここで，
各$j$について
$a_j\in \cc$
で
族$\{A_{j}\}$
は互いに交わらない可測集合とする．
同様に
\[
g(y)=\sum_{l=1}^mb_l\chi_{B_l}(y)
\]
とおく．
ここで，
各
$l$について
$b_{l}\in \cc$で，
族$\{B_{l}\}$は互いに交わらない可測集合とする．
この状況において
\[
\|f\|_p^p=\sum_{j=1}^n|a_j|^p\mu(A_j)=1
\]
と
\[
\|g\|_r^r=\sum_{l=1}^m|b_l|^r\nu(B_l)=1
\]
に注意せよ．
ここで関数$p$を
\[
\frac{1}{p(z)}=\frac{1-z}{p_0}+\frac{z}{p_1}
\]
と定義する。
さらに
$r_0$, $r_1$を$q_0$, $q_1$の共役指数として
関数$r$を
\[
\frac{1}{r(z)}=\frac{1-z}{r_0}+\frac{z}{r_1}
\]
と定義する．
そして$f_{z}$と
$g_{z}$を以下のように定義する．
\[
f_z(x)=|f(x)|^{\frac{p}{p(z)}}\frac{f(x)}{|f(x)|}
\]
\[
g_z(y)=|g(x)|^{\frac{r}{r(z)}}\frac{g(x)}{|g(x)|}. 
\]
すると，
\[
f_z(x)=\sum_{j=1}^n|a_j|^{\frac{p}{p(z)}}\frac{a_j}{|a_j|}\chi_{A_j}(x)
\]
及び
\[
g_z(y)=\sum_{j=l1}^m|b_l|^{\frac{r}{r(z)}}\frac{b_l}{|b_l|}\chi_{B_l}(y)
\]
と表すことができる．
そして
\[
I(z)=\int_{Y}g_zTf_z\ d\nu
\]
と定義する．
すると
\[
I(z)=\sum_{j=1}^n\sum_{l=1}^m|a_j|^{\frac{p}{p(z)}}\frac{a_j}{|a_j|}|b_l|^{\frac{r}{r(z)}}\frac{b_l}{|b_l|}
\]
が成り立つ．
もちろんこれは
$\{z\in \zz\mid 0\le \re z\le 1\}$
上で有界かつ連続で，その内部で正則関数である．
$c\in [0,1]$に対して
\[
M(c)=\sup_{y\in \rr}|I(c+iy)|
\]
と定義すれば，三線定理から
\[
M(c)\le M(0)^{1-c}M(1)^c
\]
である．
$M(0)$と$M(1)$の評価をしていこう．
任意の$y\in \rr$に対して
\[
|I(iy)|=\left|\int_Yg_{iy}Tf_{iy}\ d\nu\right|\le \|Tf_{iy}\|_{q_0}\|g\|_{r_0}
\le \|T\|_{p_0, q_0}\|f_{iy}\|_{p_0}\|g_{iy}\|_{r_0}
\]
を得る．
ここで
\[
\|f_{iy}\|_{p_0}^{p_0}=\sum_{j=1}^n\left||a_j|^{\frac{p}{p(iy)}}\right|^{p_0}\mu(A_j)=\sum_{j=1}^n|a_j|^{\frac{p}{p(iy)}p_0}\mu(A_j)=1
\]
\[
\|g_{iy}\|_{r_0}^{r_0}=\sum_{l=1}^m\left||b_l|^{\frac{r}{r(iy)}}\right|^{r_0}\nu(B_l)=\sum_{l=1}^n|b_l|^{\frac{r}{r(iy)}r_0}\nu(B_l)=1
\]
であるから，
\[
M(0)\le \|T\|_{p_0,q_0}
\]
を得る．
同様にして．
\[
M(1)\le \|T\|_{p_1, q_1}
\]
を得る．
よって
\[
\left|\int_{Y}gTf\ d\nu \right|
\le \|T\|_{p_0, q_0}^{1-\theta}\|T\|_{p_1.q_1}^{\theta}
\]
を得る．
$p_0=q_0=\infty$などの場合は、$f_z=f$と定義すれば良い．
\end{proof}

これの系として以下の定理がある。
\begin{thm}[Riesz--Thorin]
$(X,\mathcal{A},\mu)$と$(Y,\mathcal{B}, \nu)$を測度空間とする。
そして$p_0, q_0, p_1,q_1 \in [1,\infty] $で
$\theta\in [0,1]$とする。
この時
\[
\frac{1}{p}=\frac{1-\theta}{p_0}+\frac{\theta}{p_1}
\]
\[
\frac{1}{q}=\frac{1-\theta}{q_0}+\frac{\theta}{q_1}
\]
とする。
このとき線形写像
\[
T:L^0(X,\mu)\to L^0(Y,\nu)
\]
について
\[
\|T\|_{p,q}\le \|T\|_{p_0, q_0}^{1-\theta}\|T\|_{p_1.q_1}^{\theta}
\]
が成り立つ。
\end{thm}




\section{クラークソンの不等式と一様凸性}

\begin{thm}[Clarkson's inequality]
$p\in (1, 2]$で $1/p+1/q=1$とする。このとき
任意の$x,y\in \cc$に対して
\[
(|x+y|^q+|x-y|^q)^{\frac{1}{q}}\le 2^{\frac{1}{q}}(|x|^p+|y|^p)^{\frac{1}{p}}
\]
が成り立つ．
\end{thm}
$2$次正方行列
$A$を
\begin{proof}
\[
A = \left(
    \begin{array}{cc}
      1 & 1  \\
      1 & -1 
    \end{array}
  \right)
\]
と定義する．
すると
\[
\|A(x,y)\|_2^2=|x+y|^2+|x-y|^2=2(|x|^2+|y|^2)=2\|(x,y)\|_2^2
\]
なので
\[
\|A\|_{2,2}=2^{1/2}
\]
であり，また
\[
\|A(x,y)\|_{\infty}\le |x|+|y|=\|(x,y)\|_1
\]
であり，
\[
\|A(1,0)\|_{\infty}=\|(1,0)\|
\]
なので
\[
\|A\|_{\infty,1}=1
\]
である．
ここで，$p\in (1,2]$より，
\[
\frac{1}{p}=\frac{1-\theta}{2}+\frac{\theta}{1}
\]
となる$\theta$が存在する．
よって
\[
\frac{1}{q}=\frac{1-\theta}{2}+\frac{\theta}{\infty}
\]
である．
よって，
\[
\|A\|_{q,p}\le (2^{1/2})^{1-\theta}1^{\theta}
=2^{(1-\theta)/2}=2^{1/q}
\]
となる．
これは求める不等式と同値である．
\end{proof}


\section{以下未完成}
\begin{prop}
$a, b$を非負の実数とする．このとき
$(a^p+b^p)^{\frac{1}{p}}\le 2^{\frac{2-p}{p}}(a^q+b^q)^{\frac{1}{q}}$
\end{prop}
が成り立つ．
\begin{proof}
ヘルダーの不等式から従う。
\end{proof}

\begin{prop}[逆ミンコフスキーの不等式]
$k\in (0,1)$とする。このとき
$f, g\in L^{k}(X)$に対して
\[
\left(\int_X|f+g|^k d\ \mu\right)^{\frac{1}{k}}>\left(\int_X|f|^k \ d\mu\right)^{\frac{1}{k}}+\left(\int_X|g|^k \ d\mu\right)^{\frac{1}{k}}
\]
が成り立つ．
\end{prop}

\begin{thm}[Clarkson's inequality]
$p\in (1, 2]$で $1/p+1/q=1$とする。このとき$f,g\in L^p$に対して
\[
(\|f+g\|_p^q+\|f-g\|_p^q)^{\frac{1}{q}}\le 2^{\frac{1}{q}}(\|f\|_p^p+\|g\|_p^p)^{\frac{1}{p}}
\]
が成り立つ．
\end{thm}
\begin{proof}
$\frac{p}{q}<1$に注意しよう。
\[
\|f+g\|_p^q+\|f-g\|_p^q=\left(\int_X(|f+g|^q)^{\frac{p}{q}} \ d\mu\right)^{\frac{q}{p}}+\left(\int_X(|f-g|^q)^{\frac{p}{q}} \ d\mu\right)^{\frac{p}{q}}
\]
よって
\[
\left(\int_X(|f+g|^q)^{\frac{p}{q}} \ d\mu\right)^{\frac{q}{p}}+
\left(\int_X(|f-g|^q)^{\frac{p}{q}} \ d\mu\right)^{\frac{p}{q}}
\le 
\left(\int_X(|f+g|^q+|f-g|^q)^{\frac{p}{q}} \ d\mu\right)^{\frac{q}{p}}
\]
である。
ここで、
\[
(|f+g|^q+|f-g|^q)^{\frac{p}{q}} \le 2^{\frac{p}{q}}(|f|^p+|g|^p)
\]
なので
\[
\left(\int_X(|f+g|^q+|f-g|^q)^{\frac{p}{q}} \ d\mu\right)^{\frac{q}{p}}
\le 2^{\frac{p}{q}}\int_X|f|^p\ d\mu+\int_X|g|^p\ d\mu
=2^{\frac{p}{q}}(\|f\|_p^p+\|g\|_p^p)
\]
よって最初に挙げた不等式は成り立つ。
\end{proof}


\begin{thm}
$2\le p$で $1/p+1/q=1$とする。このとき$f,g\in L^p$に対して
\[
(\|f+g\|_p^p+\|f-g\|_p^p)^{\frac{1}{p}}\le 2^{\frac{1}{p}}(\|f\|_p^q+\|g\|_p^q)^{\frac{1}{q}}
\]
が成り立つ．
\end{thm}
\begin{proof}
$\frac{q}{p}<1$に注意しよう。

\[
\|f+g\|_p^p+\|f-g\|_p^p=\int_X|f+g|^p\ d\mu+\int_X|f-g|^p\ d\mu=
\int_X (|f+g|^p+|f-g|^p)\ d\mu
\]
ここで
\[
|f+g|^p+|f-g|^p\le 2(|f|^q+|g|^q)^{\frac{p}{q}}
\]
なので
\[
\int_X (|f+g|^p+|f-g|^p)\ d\mu
\le 2\int_X(|f|^q+|g|^q)^{\frac{p}{q}}\ d\mu
\]
ミンコフスキーに不等式から
\[
\left(\int_X(|f|^q+|g|^q)^{\frac{p}{q}}\ d\mu\right)^{\frac{q}{p}} \le
\left(\int_X|f|^p\ d\mu\right)^{\frac{q}{p}}+\left(\int_X|g|^p\ d\mu\right)^{\frac{q}{p}}=\|f\|_p^q+\|g\|_p^q
\]
よって
\[
\|f+g\|_p^p+\|f-g\|_p^p\le 2(\|f\|_p^q+\|g\|_p^q)^{\frac{p}{q}}
\]
を得る。これは先に挙げた不等式に等しい。
\end{proof}

\begin{cor}
$L_p$は一様凸である。
\end{cor}

\section{追記}

1936年に
クラークソン
\cite{MR1501880}
はクラークソンの不等式を用いて
$L^{p}$空間や
$\ell^{p}$の一様凸性を証明した．

1940年にBoas
\cite{MR0001456}
は``線形写像のリースの凸性定理''
を用いてクラークソンの定理を証明しているが，
多分これは今ふうにいうとリースソリンのことだと思う．

1956年に
Henner
\cite{MR0077087}
はクラークソンの不等式の拡張を，
なんか頑張って不等式評価して証明している．


1959年に
Morawetz
\cite{MR0256149}
は
$L^{p}$空間の一様凸性の証明に使うことのできる
新しい不等式を証明している．


1976年に
Ramaswamy
\cite{MR0492135}
は実数に関するクラークソンの不等式を用いるのではなく，
ラーデマッヒャー関数を用いて積分のクラークソンの不等式を
直接証明している．


%READ!
%myplain is the same to amsplain style. 
\nocite{*}
\bibliographystyle{myplaindoidoi}
\bibliography{cl.bib}



	

			\end{document}



\begin{comment}
[集合のやつは$\mid$を使う。]
[真なる包含関係]
\subsetneqq
[可換図式]
\[\xymatrix{
&   & (2) \ar@{=>}[rd]&  &  &  & (5) \ar@{=>}[rd]&\\ 
& (1)\ar@{=>}[ru] \ar@{=>}[rd]&  &(4)\ar@{=>}[r]&(7)& (1)\ar@{=>}[ru] \ar@{=>}[rd]& & (7)&\\ 
&   & (3) \ar@{=>}[ur]&  &  &  &(6) \ar@{=>}[ur]&
}\]

[\bigin \endを使わない方法]

{\prop[aa] aaa}
で
\begin{prop}[aa]
aaa
\end{prop}
と同じ\df \thm \proof でも同様

[直和]
$\amalg$

[同値関係でよく使う波線]
\sim

[引数付きの新命令]
\newcommand{\prn}[1]{\|#1\|_p}

[番号のリセット]

\setcounter{equation}{0}


[字体の変更]

{\rm hogehoge}と使う。

\rm ローマン
\it イタリック
\sl 斜体

[supp]

\newcommand{\supp}{\text{{\rm supp}}}

[中央揃えの改行]

	\begin{eqnarray*}
		hoge1&=&hogehoge
		hoge2&=&hogehofe

	\end{eqnarray*}

&&の部分で揃えられる。


[\tag]
\begin{equation}
f(x) = ax^2 + bx + c \tag{1.2.3}
\end{equation}



[場合分け]

	\begin{cases}
		hoge1 & \text{状況１}\\
		hoge2 & \text{状況２}\\
		・
		・
		・
	\end{cases}



\end{comment}





\begin{enumerate}
  \item 
\end{enumerate}
